\subsection{Compose methods, state controls, and event intent in one authored program}
\label{sec:common-program-authoring}

The preceding examples established four separate boundaries: a method can be
written in complete text, a registry can retain its source and identity, a
state invariant can reject a bad book, and a selected event contract can retain
the caller's full intent. A usable program must express those choices together.
Otherwise a typed fixture may pass while the proposed language has no way to
write the same program or to show which semantic owner supplied a missing field.

\code{AuthoredCompanyText/1} closes this composition boundary with a bounded
source format. The program imports the opening book, report schemas, statement
mapping and memory definition by exact SHA-256 identities. It selects one
\code{BookStateInvariant/1} resource explicitly. These imports are facts and
checked rule resources; they do not silently generate methods or claim that a
small synthetic state is a complete disclosure history.

Each registry entry contains the complete
\code{CompanyMethodAuthoring/1} source block. The entry records its lowered
method digest, source digest, transition profile, input declarations, account
classes, clock and separate entry-schema resources. A registry version therefore
cannot hide a transition change. The method header supplies the transition
profile, while the program author chooses the registry and method for each
event. Entry schemas remain distinct from the shared report schemas because
combining them would change the identity of a report request.

For the selected receivable contract, the same program also declares the native
implementation identity, role-to-account map, unit-cost pointer and policy.
Every event carries both its projected method inputs and its complete original
intent. The latter retains the customer, claim and interval assertions that a
specialized method may remove from its calculation inputs. The runtime then
checks the selected native lifecycle on the current book and requires the
authored projected request to equal the declared projection of that full
request. An edited half-collection method therefore remains executable under
the guarded target, where it changes the resulting claim and cash, but it is
refused when the caller explicitly selected the full-payment native contract.
This is a contract comparison, not an inference that either method is the
correct accounting treatment.

Reports, notes, carry boundaries and memory policies reuse the existing command
forms. A reset policy can obtain its type only from an explicitly present
opening memory leaf. Unknown event inputs retain their runtime unresolved
status. Duplicate declarations, unregistered aliases, changed resources and
incomplete source blocks are errors with a stage and UTF-8 byte location.
Method source is hashed from its original bytes; the
verifier checks outer, method and declaration spans after Unicode and CRLF
roundtrips.

Two complete eleven-command programs reconstruct their saved typed programs
and runtime results exactly. The run includes nineteen structural refusals and
six matched native refusals across the two charts, direct method editing,
report-note preservation and failed-prefix checks. These finite comparisons
show that complete methods, explicit registries, selected state rules and full
event intent can share one source boundary. They do not provide general
invariant authoring, whole-company source completeness, accounting admission,
production readiness or human acceptance. The following successor supplies
explicit source for the current finite invariant family.

\begingroup\raggedright
The field inventory and grammar design are retained under
\path{data/accounting-master/common-program-authoring-audit-20261026/}.
The implementation and checks are under
\path{data/accounting-master/common-program-authoring-20261026/}.

\par
\endgroup

\subsection{Write the invariant that a company program selects}
\label{sec:invariant-source-authoring}

Selecting a rule by its digest preserves its identity, but gives the author no
syntax for changing the rule itself. The claim example makes the missing
operation concrete. Receivables are 30 and the proposed claim is 31. An author
must be able to state that all claim amounts sum to receivables, specify which
state fields supply those amounts, and choose how an unknown total affects the
program. Naming an imported algorithm alone does not express those choices.

\code{BookStateInvariantText/1} writes the existing finite invariant as
explicit declarations. Global declarations identify the book date, balances or
memory leaves. Each collection identifies a memory map or the book's seen-event
sequence. Its fields declare their types, contexts, provenance and bindings to
map keys, record fields or sequence values. Uniqueness, ordering and reference
constraints name the collections and fields they compare. Per-row blocks
contain guarded predicates and explicit row/global bindings. A final aggregate
block names sums and global inputs, then states the predicates they must satisfy.
Residual prerequisites remain attached to the row, collection, aggregate or
whole invariant at which they apply.

The frontend lowers these declarations to the same \code{BookStateInvariant/1}
definition. The retained collection, total and guarded-requirement evaluators
still determine every predicate's meaning. No callback or second arithmetic
engine is introduced. \code{AuthoredCompanyText/2} embeds that source between
\code{begin\_invariant\_source} and \code{end\_invariant\_source}, replacing
version 1's selection of an imported whole-rule resource. The compiler constructs
the exact typed rule internally and passes it to the original program compiler.
Opening books, schemas and other data resources remain explicit imports.

For the claim total, the aggregate predicate is equality between the sum of
remaining claim amounts and the receivable balance. Changing its false outcome
from \code{invalid} to \code{unsupported} changes the reason the program stops;
the proposed sale still fails and its journal never enters the accepted prefix.
A more consequential edit replaces equality with ``receivables are no greater
than the claim total.'' That weaker predicate permits the 31-versus-30 case.
The source now controls the rule's truth, but the altered rule has lost the
original accounting reconciliation. The comparison is an authoring test, not a
recommendation to weaken that control. Selecting a well-formed rule never proves
that it is the right rule for a company.

Compilation and state binding remain separate. A uniqueness rule over dates is
a valid typed definition. If the author binds those dates to actual text map
keys, the native binder refuses the mismatch, including when the map is empty.
Partial populations and unknown amounts remain unresolved; an observed subtotal
does not become the complete total. These distinctions are preserved by reusing
the original typed evaluators rather than by treating a successful parse as
evidence that the book satisfies its controls.

The two company sources reconstruct their original invariants and all 22
complete command results. All sixteen inconsistent-state examples retain their
full refusal results, unchanged opening and unexecuted suffix. Fourteen source
controls exercise duplicate, missing, unbound, incompatible and oversized
declarations. Both direct rule edits agree with their explicitly edited typed
counterparts, and the two partial/unknown cases retain their full unresolved
results. Source and nested declaration spans are checked against UTF-8 bytes
with both ordinary and CRLF line endings. Whitespace can change the source
digest while leaving the typed program unchanged.

This is a bounded authoring language for the existing state rules. It does not
add arbitrary algorithms, establish completeness of a real company's claim
population or supply accounting authority. Wider rule families and independent
implementation remain open. The grammar decision, source examples, complete
comparisons and unchanged-source checks are retained under
\path{data/accounting-master/invariant-source-authoring-20261026/}.


\subsection{Write the selected company rules and commands in one source}
\label{sec:selected-company-source}

A generated method can calculate the right journal while leaving the modeller
unable to state the rule that generated it. In the claim example, the customer
check, the source of a claim identifier and the order of assessment updates are
substantive choices. Naming a template by its digest records which choice was
made, but does not let an author write or change that choice. The selected
company source therefore includes the opening conditions, state invariant,
eligibility stages and construction recipes themselves. The earlier source
frontends remain useful predecessor comparisons; this successor targets the
declared opening and execution path together. The opening template in
this profile still requires an empty claim map. Section~\ref{sec:complete-claims-opening}
extends that boundary for explicitly complete supplied carried claims.

Consider a collection whose claim identifier is absent. Reading that claim's
customer immediately would fail for a missing row, obscuring the intended
refusal that the claim does not exist. In \code{OrderedEligibilityText/1}, the
author writes a membership stage before the customer stage. A binding such as
\code{(map\_contains "/claims" "claim\_id")} supplies the membership
predicate. A later \code{map\_field} binding names the same map and event
input, the customer field and the preceding presence stage. Compilation accepts
that lookup only when the earlier stage unconditionally requires the matching
membership fact. All stages compile, including ones that an earlier refusal
would prevent from running. An empty map cannot hide a malformed field lookup.

The full event-input types, stage identifiers, failure reasons and guarded
predicates are also declarations. The retained order first resolves the outer
event structure and unknown inputs, then approval and the operation's specific
conditions. Moving claim existence ahead of the sealed-day condition changes
which reason is returned when both fail. Negating the approval predicate
changes which events are eligible. These edits demonstrate that the source
controls the selected rule; they are different rules, not equivalent ways to
spell the retained contract. The six eligibility sources reconstruct their
complete typed templates and preserve 38 eligible and 34 refused gate results.

Opening admission has a related but distinct job. \code{OpeningAdmissionText/1}
states structural checks for the input fields, policy values, row and clock
schemas, roles and report mapping, alongside typed predicates for the dates
and amounts. The following excerpt is an actual declaration for the final
command-count condition; its position follows the other opening checks.

\begin{quote}\small\ttfamily
stage "commands" commands "command\_limit";\\
check 64;\\
end\_stage;
\end{quote}

Other structural stages declare their paths and expected values directly.
Predicate stages use the same guarded-expression grammar as eligibility.
Dependency declarations name the exact reused semantic implementations, while
identities for locally authored predicates and templates are computed from the
definitions. The two 21-stage sources preserve all 82 complete opening results.
Ten inputs are admitted, including those whose future command will later fail.
Neither successful parsing nor admitted opening conditions establish that the
whole future program succeeds.

Construction recipes express the next step: how a permitted event supplies the
values needed by a method. \code{MethodRecipeText/1} keeps the existing
method declarations for monetary calculations, postings, cash flows, guarded
conditions and memory writes, and adds explicitly typed parameters. For
example, a customer parameter can read the event's customer input, and a
memory write can refer to that parameter. No event value is interpreted as
source code. A pointer parameter combines literal path segments and previously
declared text parameters; the existing pointer encoder escapes slashes and
tildes in a claim identifier before using it as a path.

The distinction matters for the sale and collection methods. The sale writes
the customer's name, due date, event identifier and issue date into a new
claim. The collection reads the selected claim's remaining amount through a
parameterized pointer. These values are supplied by the actual event and
current book. The accounting expressions and account choices remain visible
in the recipe. An author need not reproduce a serialized method tree or
maintain its selection hash manually. Missing, unused, forward or mistyped
parameters are refused by the original contextual constructor.

An assessment requires a sequence of writes over the current claims.
\code{OrderedExpansionText/1} contains the complete base recipe, map path,
row types, per-row replacements and final replacements. Each replacement
declares its identifier, path and value expression. A row expression can read
a field, the current key, an event input or event date, and choose between
typed branches of a date condition. Both branches are checked even when only
one will be evaluated. The explicit assessment-key order determines the write
sequence; a duplicate or incomplete order refuses construction. A malformed
row expression also refuses when the map happens to be empty.

Four scalar recipe sources and two ordered-expansion sources reconstruct their
complete templates. Across 18 scalar and 20 expansion cases, the full
construction results, transition requests and transitions equal the retained
typed comparisons. These checks include identifiers containing slashes, tildes
and Unicode. Source lowering reports syntax and structure only; the current
book, event, schemas and memory definition remain necessary for contextual
construction. Thus a reusable recipe is not prematurely described as a method
valid for every possible state.

\code{SelectedCompanyText/1} joins these declarations. It imports the opening
book, schemas, statement mapping and memory definition by exact content
identity. Those four resources contain facts and types. The opening rule,
invariant, eligibility and construction algorithms are inline source blocks.
Roles, unit-cost pointer and policy define the selected company conventions.
Events retain their full typed inputs, provenance and declaration order;
reports and carry commands retain their dates, notes, equity movements and
explicit memory policies. The native-shaped admission input is derived from
these same facts and projected commands, rather than authored as an unrelated
second copy.

Lowering creates the full program and its selection without running commands.
Execution then reuses the opening-to-program binding and the existing
eligibility, invariant, construction, transition, report and carry evaluators.
For four successful programs, all 44 complete command results are preserved.
For 26 adverse programs, the stopped prefix, current and opening states,
refused request where present and unexecuted suffix are preserved as well.
All native lifecycle functions and the operating-method builder are disabled
during these comparisons. Changing the first sale's approval directly in the
company source stops the program before any command succeeds.

One provenance distinction must remain visible. Canonical content identities
ignore dictionary insertion order, whereas the outer execution binding also
records the serialized order of the complete supplied program. The source
parsers reconstruct some definition dictionaries in a different order. Their
canonical content is identical, but that one outer order digest changes. The
comparison recomputes and retains both order identities, and permits a
difference only at \code{program\_order\_sha256}. Every other result field
matches, including the separate digest of event-input name order, which must
remain unchanged because it affects the order of reported unknown inputs.
This is a precise provenance exception, not a relaxation of financial-result
or failure comparison.

Source identity retains the original UTF-8 bytes. Whitespace and CRLF line
endings can change that identity while leaving the lowered program and its
results equal. Errors inside a nested rule or recipe identify the corresponding
bytes in the full company source. A lexical error can identify its token; a
semantic error that is known only at stage or document scope is reported at
that scope. The corrected nested-invariant reader described in
Section~\ref{sec:complete-claims-opening} also preserves quoted Unicode
separator characters as data rather than treating them as source line endings. The checked sources reject duplicate or missing declarations,
changed or unused imported resources and incompatible operation selections.

The comparison supports this finite synthetic company language. It does not
make an edited rule accounting-equivalent to the retained rule, supply missing
facts about a real company's claim population, or complete the broader fiscal,
sector and measurement algorithms. Independent implementability and reader
acceptance remain open. The complete source examples and result pairs are
retained under
\begingroup\raggedright
\path{data/accounting-master/selected-company-source-20261026/}.
The component evidence is under
\path{data/accounting-master/ordered-eligibility-source-20261026/},
\path{data/accounting-master/opening-admission-source-20261026/}
and \path{data/accounting-master/method-recipe-source-20261026/}.
\par\endgroup


\subsection{Begin a reporting period with live claims}
\label{sec:complete-claims-opening}

A new reporting period need not begin with an empty receivables register.
Customers may still owe amounts from sales in the preceding year. The closing
procedure carries those claims forward with their identifiers, issue dates,
remaining amounts and assessment history. An opening rule that requires an
empty register would reject this ordinary continuation even when the carried
book and opening statements are internally consistent.

The retained examples make that conflict concrete. Two company charts carry
one live claim with a remaining amount of 20 into 1 January 2028. Two further
programs carry three claims whose remaining amounts are 13 each, giving a
receivables balance of 39. The claims were issued during 2027 and assessed on
31 December 2027. The new reporting period ends exclusively on 1 January 2029;
its period flows are empty, while the claim records and previously seen event
identities remain. The original continuation contains three commands in each
one-claim case and two in each three-claim case. Those ten commands supply an
exact comparison for the opening extension.

The selected empty-opening template rejects all four books at
\code{empty\_claims}. Removing that stage alone exposes a second empty-specific
condition: \code{claim\_total} requires the receivables balance to be zero.
Removing both conditions makes the remaining opening checks pass, but it also
admits five deliberately corrupted books. The corruptions are a wrong
receivables total, a negative remaining claim, a missing sale identity, an
issue date after the opening and an inconsistent overdue amount. The full
selected state invariant rejects each corruption. Thus the required change
is not simply to permit a nonempty map. It is to retain the relationship
between every supplied claim and the book as a whole.

\code{CompleteClaimsOpening/1} makes that change under an independently
selected contract. The selection identifies the original empty-opening
template and the full state invariant. The new contract carries both
definitions, their canonical identities, the exact two excluded stage names
and an explicit description of its scope. Before deriving the revised
structural template, the adapter checks that \code{empty\_claims} is the
reviewed equality with the empty claim map and that \code{claim\_total} is
the reviewed zero-receivables predicate. Every other stage remains unchanged.
A newly supplied predicate does not acquire reviewed status merely because
the caller computes its digest.

The accompanying \code{SuppliedClaimPopulation/1} declaration states whether
the supplied internal claim map is complete, partial or unknown. It names
\code{/claims}, identifies the exact opening book and supplies provenance.
Only an explicit \code{complete} declaration proceeds under this profile;
partial and unknown declarations return unresolved before any command runs.
Completeness here concerns the supplied internal state. A populated map, a
balanced book and a successful continuation cannot establish that a real
issuer has disclosed every claim. That different question still belongs to
source-constrained reconstruction.

Admission first runs the retained structural and opening-report checks on the
same input. It then evaluates the full selected invariant over the unchanged
book, schemas and memory definition. The invariant checks the aggregate
receivables relationship and the claim conditions that the simple relaxation
missed. A malformed schema or incomplete statement mapping therefore still
fails its original check. A mid-period snapshot also remains unsuitable as a
new-period opening even if its claims satisfy the invariant. This profile
deliberately checks structure and the opening report before the full
invariant; it does not promise an older profile's failure precedence when
several conditions fail together.

The program wrapper, \code{AdmittedCompleteClaimsProgram/1}, retains the
complete admission result. After admission succeeds, it delegates to the
existing opening-to-program binding using the derived structural template
and its corresponding internal selection. That binding again checks shared
facts, schemas, mapping, memory definition, roles, policy, rules, invariant,
complete projected commands and event-input order before the unchanged
composer runs. A caller-supplied success flag cannot replace these checks.
The retained delegated result exposes the requests, transitions, reports,
carry state, trace and any stopped suffix.

The same distinction is explicit in company source.
\code{CompleteClaimsCompanyText/1} surrounds the full selected company
source with a reviewed-selection identity, population declaration and contract
residuals. The declaration below supplies the population premise in a checked
carried example; the opening-book identity is computed from the exact imported
book rather than typed a second time.

\begin{quote}\small\ttfamily
population "SuppliedClaimPopulation/1" complete "/claims"\\
\hspace*{1em}"Exact retained synthetic carry; complete supplied internal map";
\end{quote}

This display wraps one source declaration for readability. The complete
program begins with the text and target profiles, then names the expected
reviewed-selection digest. The caller supplies that reviewed selection
separately. Between \code{begin\_selected\_company} and
\code{end\_selected\_company}, the author writes the complete opening,
invariant, eligibility, recipe and command source described in
Section~\ref{sec:selected-company-source}. Its only imported resources remain
the opening book, schemas, statement mapping and memory definition. An extra
resource containing an opaque algorithm is refused.

Lowering checks that both the named identity and the nested source selection
match the independently supplied selection. Editing an opening condition or
invariant cannot silently replace that reviewed choice. Lowering also verifies
the structural contract, but does not admit the opening or run commands.
Execution invokes admission anew. UTF-8 bytes and line endings remain part of
source identity; nested errors retain their original parent-source locations.
The invariant parser supplies declaration locations, while the opening lexer
can identify an offending token. The outer source preserves that difference
instead of asserting precision the inner parser does not provide.

A source line ends at LF, with an optional preceding CR; a quoted string may
contain other Unicode separator characters. The inherited invariant parser
incorrectly split such strings at U+0085, U+2028 and U+2029. The corrected
reader and nested rule renderer keep these characters inside the string.
Nine comparisons place each character in an invariant residual, a nested
rule residual and an input source reference. Each reproduces the earlier
parser failure and preserves the corrected typed program and complete result.
All eight retained programs and 32 prior controls also remain exact. This
repair changes source handling only; the financial evaluators are unchanged.
The repaired invariant and complete company source implementations are under
\begingroup\raggedright
\path{data/accounting-master/invariant-unicode-repair-20261026/}.
\par\endgroup

Four carried and four empty-opening source programs reproduce the complete
typed results. The carried programs retain all ten continuation commands;
the empty-opening programs retain all 44 commands. The five corrupted states
still refuse before execution. Native lifecycle callbacks and the operating
method builder are disabled throughout these comparisons. The result
comparison permits differences only in the outer and delegated
\code{program\_order\_sha256} fields: both record the full program's
serialization order, which source reconstruction can change while preserving
canonical content. Both identities are independently recomputed and retained.
Every other field, including complete execution and event-input-order identity,
matches the typed comparison.

The checked successor therefore admits these complete supplied carried books
without weakening the selected invariant. It leaves partial actual histories,
arbitrary alternative financial rules and accounting applicability unresolved.
The carried-state audit, typed contracts and full source comparisons are kept
together with their explicit refusal cases under
\begingroup\raggedright
\path{data/accounting-master/nonempty-opening-audit-20261026/},
\path{data/accounting-master/complete-claims-opening-20261026/}
and \path{data/accounting-master/complete-claims-source-20261026/}.
\par\endgroup


\subsection{Keep each assessment attached to its historical terms}
\label{sec:nested-term-invariants}

A claim can remain outstanding while its due date changes. Checking its
current balance and current due date is therefore insufficient to explain an
earlier overdue assessment. In the retained synthetic example, a claim issued
on 2 January 2027 is initially due on 31 January. An imported assessment dated
28 February records 20 remaining and 20 overdue. Prospective extensions become
effective on 10 March and 10 April, moving the due date first to 30 April and
then to 31 May. Applying the current 31 May due date to the February assessment
would incorrectly turn its overdue amount into zero. The historical assessment
must remain attached to the initial term.

The term history is stored inside the claim. Each term has its own local
identifier, effective date, due date, predecessor and event identity. Each
assessment has a date, referenced term, remaining amount, overdue amount,
origin and event identity. Local identifiers alone do not identify these
records: another claim may use the same identifier. Nor does insertion order
establish chronology. A declaration must name the parent claim and select a
term using its effective date before comparing the financial amounts.

\paragraph{Expose the supplied rows without choosing a history.}
\code{NestedMapProjection/1} binds a declared parent map and nested child map
to typed scalar rows. Bindings may read the parent key, child key, a parent
field or a child field. The accompanying lineage retains both raw keys and
escaped source paths. Each parent also has a group record with its child-map
path, first row and row count. An empty child map still produces a group;
an empty parent map produces no groups. This distinction lets a rule require
terms for every supplied claim even when one claim has no terms.

Projection preserves supplied insertion order and makes no chronological
selection. Declarations are checked against the memory definition even when
there are no rows. An explicit population declaration identifies the exact
book and both map paths; partial or unknown coverage remains unresolved.
Complete supplied internal maps are sufficient for these checks, but do not
establish that an issuer's disclosures reveal all contracts or past events.

\paragraph{Select within one parent and one declared boundary.}
\code{GroupedRowQuery/1} provides count, typed match and greatest-row queries.
The query names the parent key, escaped parent path and child-map path; all
must agree. Count includes all rows in that group. Match compares one declared
field with one value of exactly the same type. Greatest orders dates or exact
rational values, optionally below an explicitly inclusive or exclusive upper
bound. Units, scale and context must agree before a rational comparison.

The February assessment selects the greatest effective date no later than
28 February, which is 2 January. A predecessor query instead uses a strict
upper bound: the predecessor of the term effective on 10 April is the term
effective on 10 March. A query without an upper bound identifies the latest
term. None of these operations sorts or replaces the stored source order.
Candidate and winner indices retain that order, and a unique result includes
the whole selected row and its source lineage.

Equal greatest dates produce an ambiguous result. Duplicate matches do the
same; insertion order is never a tie breaker. A group absent from a complete
supplied population is invalid. A present group with no qualifying row returns
explicit \code{none} when the declaration permits it, or invalid when the row
is required. An explicit unknown input remains unresolved; omitting a required
input is a malformed declaration. Every evaluation recomputes its projection
from the bound book, schemas, types and coverage. A caller cannot substitute
previously computed rows after changing the source.

\paragraph{Use absence without treating uncertainty as absence.}
\code{QueryGuardedInvariant/1} binds query results into the preserved guarded
requirement evaluator. Count becomes an exact dimensionless rational. A unique
row supplies \code{has\_row=true} and known selected fields. Allowed
\code{none} supplies \code{has\_row=false}, while its selected fields remain
unknown. An ambiguous or unresolved query supplies unknown values, including
an unknown presence indicator. Thus a tie cannot masquerade as an empty group.
The declaration explicitly chooses whether ambiguity is invalid or unresolved.

A false guard makes a predicate inactive; it does not prove that predicate.
The separate \code{required\_known} declaration can require a selected field
even when its predicate is inactive. Invalid queries fail structurally before
the predicates run, and unresolved queries keep the overall result unresolved
even under a false guard. Allowed absence is different: it can legitimately
activate an empty-group rule while leaving a selected-row rule inactive. The
result retains the original guard findings and distinguishes
\code{no\_active\_requirements} from satisfied requirements.

\paragraph{Declare the rule once for every supplied claim and assessment.}
Finite queries with handwritten claim identities would miss newly supplied
records. \code{NestedQueryInvariantTemplate/1} instead iterates over all groups
or all child rows of a declared complete projection. Typed captures read the
current parent or child identity, scalar fields and explicit book dates or
memory values. Every query's full group identity is bound to that parent.
Scalar substitutions can fill only a typed match value, typed bound or the
value of a typed constant already present in a guarded predicate. They cannot
replace operators, types, bindings or failure consequences. Duplicate or
unsafe substitutions refuse before iteration, including on empty populations.

Each instance retains its captured values, source paths, resulting declaration
and complete query and predicate results. The reviewed selection identifies
the template as a whole. The instance selection is explicitly derived by the
checked substitution procedure; it is not described as an independent review
of every resulting instance. The initial checked profile permits at most
64 instances and refuses overflow without a partially successful result.
This bound is an implementation choice for the finite probe, not a restriction
on the financial concepts that the product must eventually represent.

The term-row declaration uses a unique match at each effective date to reject
ties. It selects the greatest strictly earlier term and checks the stored
predecessor against that identity; when no predecessor exists, the explicit
initial sentinel must be \code{none}. It also checks issue, effective, due and
book dates. The assessment-row declaration checks uniqueness of its date,
selects the greatest applicable term, preserves historical overdue, requires
nonnegative assessed remaining amounts and distinguishes imported snapshots
from recorded events. Parent-group declarations check current terms, latest
assessments, assessed flags, summary amounts and empty-assessment sentinels.

Event membership uses \code{ProjectedStateCollections/1} to connect nested
rows to the existing collection-reference evaluator. It binds complete term
rows and the book's
seen-event sequence. An explicitly typed equality filter selects assessments
whose origin is \code{recorded\_event}; their event identities must be seen.
Imported snapshots have a separately declared unknown-event condition. A
filtered collection is complete only relative to its complete supplied source
and explicit predicate. The filter does not turn a partial source into a
complete collection, and its types are checked even when it selects no rows.

\paragraph{Join the six semantic checks.}
\code{TermStateSemanticBundle/1} joins six reviewed components: common claim
conditions and the receivables total, the collection-clock bound, term history,
historical assessments, parent summaries and event membership. The clock reads
the actual exclusive fiscal endpoint rather than assuming a calendar year.
The common claim declaration retains nonnegative remaining amounts, sale-event
membership, issue dates and the exact receivables sum. Historical assessment
rules own overdue; the earlier scalar rule based on current due is not applied
to these amended-term histories.

All 22 retained states satisfy this semantic bundle. Fifty-one complete
comparisons preserve every subassessment, and 27 targeted semantic corruptions
are rejected by both the independently run native validator and the declared
checks. The candidate evaluations disable native financial callbacks. A
separate discriminator changes the native migration premise: native validation
rejects that premise while the book's declared semantic relations still hold.
This difference identifies a separate structural condition rather than an
alternative form of opening admission. The semantic bundle by itself grants
no execution or accounting admission. Section~\ref{sec:selected-term-execution}
supplies the selected structural contract, direct nested source and bounded
event/carry composition while retaining actual-source and accounting limits.

Complete declarations, original-state comparisons, structural controls and
the native-oracle outcomes are retained under
\begingroup\raggedright
\path{data/accounting-master/grouped-row-query-20261026/},
\path{data/accounting-master/query-guard-composition-20261026/},
\path{data/accounting-master/nested-query-template-20261026/},
\path{data/accounting-master/nested-obligation-completion-20261026/}
and \path{data/accounting-master/term-semantic-composition-20261026/}.
\par\endgroup


\subsection{Carry the historical terms through a complete command}
\label{sec:selected-term-execution}

The historical checks in Section~\ref{sec:nested-term-invariants} answer whether
the supplied terms and assessments fit together. They do not, by themselves,
justify changing the claim. An amendment also needs a valid current state,
an identified customer and prior term, an admissible effective date, and a
complete description of the writes it will perform. A collection must change
the remaining receivable without rewriting its earlier assessments. At a
period boundary, closing entries and the next assessment window must agree.
These are connected requirements: accepting each snapshot in isolation would
leave the correspondence between the event and the financial history unchecked.

The selected continuation described here joins those checks for the retained
synthetic claim. It covers prospective due-date amendments, collections,
assessments and period carry. The accounting premise is still deliberately
narrow: an amendment changes the supplied terms without recognising a
modification gain or loss, interest or impairment. Actual contractual
applicability and the appropriate accounting treatment require their own
evidence. Within this premise, the implementation question is precise: can
explicit declarations produce the complete retained requests and results while
the native financial builders are disabled?

\paragraph{Establish which supplied state the declaration accepts.}
The selected state contract combines structural checks with the six historical
and financial checks already described. Structural checks require the exact
state fields and profile, a typed book, the declared term-claim and clock
schemas, the operating account roles, a positive cost leaf with the expected
units, and the stated migration fields and premise. The account selection is
explicit. A different cost pointer can be valid under the native validator
and still fall outside this selected contract; refusal then means that the
selection does not cover that configuration.

Existing admission stages evaluate most of these conditions on a normalized
view in which the original \code{book} field is named \code{opening}. The
original envelope is checked first, and normalization preserves every value
and its insertion order. The reused stage engine's internal success label does
not turn a mid-period snapshot into a new-period opening. The enclosing result
reports satisfaction of the selected snapshot contract and retains false
opening, source and accounting admission flags.

Two additional conditions are explicit. The migration source-book identifier
must have the syntax of a SHA-256 digest, and every term premise must contain
nonblank text. The distinction between text and nonblank text matters: a
space-only premise can satisfy the generic text type while failing the native
premise requirement. Checking a digest's syntax establishes neither the
existence nor the truth of the referenced source. Similarly, complete supplied
claim and nested-row populations remain a stated input scope. A partial
population leaves the command unresolved before construction.

\paragraph{Resolve the event against the same history.}
For an amendment, eligibility first checks the event's identity and date, the
complete typed inputs, and the supplied approval. Unknown inputs produce an
unresolved result. An explicitly unselected approval produces an unsupported
result. These outcomes differ from invalid state, stale selection, duplicate
event or contradictory dates.

The event must name an existing claim belonging to the stated customer. Its
prior-term input must equal the claim's current term. The new term identifier
must be nonblank, unused and different from the sentinel \code{none}; the
premise must also be nonblank. The new due date cannot precede the event's
effective date. Finally, the effective date must be strictly later than that
of the selected prior term. Equal dates would make the prospective sequence
ambiguous and are refused, even when the supplied snapshot is otherwise valid.

A nested lookup is permitted only after the declaration establishes the
presence of the parent claim and the selected child term. This is an ordering
requirement on the checks, not an inference from a successful dictionary
lookup. The following binding fragment names the prior term's effective date;
the two referenced stages must appear earlier and establish the corresponding
memberships. The surrounding rule declares its inputs and comparison.

\begin{verbatim}
binding "previous";
  kind nested_map_field;
  pointer "/claims";
  key_input "claim_id";
  parent_after_stage "claim_exists";
  child_pointer "/terms";
  child_key_input "prior_term";
  field "effective";
  after_stage "prior_exists";
end;
\end{verbatim}

The selected eligibility engine checks those dependencies when compiling the
rule. It then resolves the bound values against the actual book and event.
The preserved guarded-predicate evaluator owns the comparisons. Adding nested
membership and field lookup does not introduce a second arithmetic or date
language.

\paragraph{Construct the complete nested change.}
An eligible amendment constructs three ordered memory operations: insert the
new term record, replace the claim's current-term identifier, and replace its
current due date. The inserted record states due date, effective date, event
identifier, predecessor and premise. Each value comes from an explicitly typed
parameter. The parameterized constructor builds the full method and typed
transition request; it does not import a stored native method as the answer
for the current event.

An assessment has a different construction. For every supplied claim, it
inserts a complete assessment record and replaces five summary fields:
latest assessment, assessed flag, assessment date, assessed remaining amount
and overdue amount. The record retains the applicable term identifier, so
later amendments cannot change the basis of this assessment. A final operation
advances the assessment cursor to the end of the checked window. In the
retained one-claim case this produces seven writes. Two claims produce thirteen
writes, with a single sequence of operation identifiers across both claims and
the suffix.

The previous ordered expander admitted only scalar fields in each row and
scalar replacements. It could not represent an insertion into the nested
assessment map. Its successor therefore checks a declared nested row schema
and recursively checks complete write templates. A record path must name a
declared field; a dynamic map key must be typed text. Insertions target map
members and supply the entire declared record. Replacements target scalar
leaves with the correct type. Both conditional branches are checked, including
the branch not taken by the current claim.

Validation precedes expansion even when the claim map is empty. An empty map
then emits only the cursor suffix, while an invalid nested record template
still fails. The existing 256-write limit is retained: a construction that
would emit 257 writes refuses as a whole. This limit is an implementation
boundary, not a conclusion that a larger company's history is economically
impossible. Claim iteration order is explicit and is preserved in the writes;
it does not determine which term was effective on an earlier date.

\paragraph{Check the proposed result before continuing.}
The command composition checks the selected snapshot before eligibility and
again after the typed transition. A well-typed but incorrect predecessor or
overdue write can pass construction and still fail the resulting historical
requirements. In that case the result retains diagnostic transition evidence
but does not return an accepted top-level after-state. The caller cannot
continue the sequence from that rejected proposal. The successful path remains
a provisional calculation; it performs no durable commit and does not establish
source or accounting admission.

Period carry follows the same discipline. Its report request must bind to the
exact closing book. The selected policy preserves all typed memory leaves,
including immutable term and assessment histories, and resets only the
assessment cursor to the next period's start. The preserved report and carry
engines produce the closing and opening books, which are checked against the
selected snapshot contract. Completeness of the supplied population is carried
forward only from the accepted pre-state through this declared transition;
it is not manufactured for an unrelated supplied book.

Chronology is visible in the complete comparison. Each company's retained
sequence contains nine events, with carry between the seventh and eighth
events. The continuation uses the computed output of each command as the next
input, including carry. Appending carry after all nine events, or resetting
the input to a stored snapshot before each event, would not verify this chain.
The final state agrees only after the commands are placed at their actual
positions across the two fiscal periods.

\paragraph{Make the selected behavior directly authorable.}
The state-contract source spells out structural stages, nested projections,
queries, captures, template substitutions, guarded requirements and collection
constraints. Event-rule source separately spells out eligibility bindings,
typed parameters, method conditions, monetary effects and nested writes.
Carry policy has its own short declaration. These source forms contain the
algorithm-bearing declarations rather than an opaque import of a completed
financial algorithm. JSON tokens remain available for explicit types, contexts,
schema expectations and scalar data.

The frontend preserves authored declaration order and exact Unicode strings.
Only line-feed characters split source lines; carriage-return/line-feed input
is also accepted without mistaking a Unicode line separator inside a string
for a new declaration. Locations record physical line numbers and UTF-8 byte
slices. Lexical errors identify tokens or declarations; failures reported by
an opaque contextual validator identify the responsible source component,
without pretending to identify a more precise field.

Parsing, eligibility compilation and contextual construction remain distinct
states. A parsed construction is not reported as contextually compiled until
its event and book have been supplied. Execution also requires the independently
supplied selection of the combined contract and rules. Changing a source file
does not authorize its changed declaration: the old selection refuses it.
Source diagnostics are retained alongside the complete typed execution result,
so adding location information does not obscure the comparison with that result.

The finite evidence establishes this selected continuation. Structural checks
and direct state-contract source each reproduce 74 complete comparisons. The
nested constructor reproduces ten retained assessment requests and transitions.
Full command composition and its source counterpart each reproduce two
continuous traces, comprising eighteen events and two carries. Malformed
sources, incorrect bindings, partial populations, stale selections, nested
write errors and the equal-effective-date boundary have separate controls.
Native financial callbacks are disabled during candidate execution. Shared
parser and typed-engine defects remain possible; this is bounded conformance
evidence, not independent acceptance or whole-company coverage.

\begingroup\raggedright
The complete sources, expected declarations, failures and results are retained
under \path{data/accounting-master/term-structural-admission-20261026/},
\path{data/accounting-master/nested-contract-source-20261026/},
\path{data/accounting-master/nested-assessment-construction-20261026/},
\path{data/accounting-master/term-full-event-binding-20261026/} and
\path{data/accounting-master/term-event-source-20261026/}.
\par\endgroup


\subsection{Preserve payment history when financing and remedies change}
\label{sec:selected-financing-history}

A company that pays late can end with the same cash, debt and equity as a
company that pays on time. The ending balance sheet does not distinguish the
two histories. A simulator must retain what was due, what remained unpaid
when payment was missed, what was subsequently paid, and which contractual
consequences were actually established. Otherwise a successful funding decision
can silently make a past failure disappear.

Consider the supplied zero-coupon promise used in the financing comparison.
The opening book contains cash of 60, inventory of 100, debt with a carrying
amount of 80 and equity of 80. Contractual face is 100, due on 30 June 2027.
A single explicit measurement interval increases carrying debt and interest
expense by 20 at maturity. Forty new shares are subscribed for 40 in cash.
There are no fees, taxes, penalties or modification effects in this example;
the price and funding date are supplied scenario assumptions. These choices
make payment timing visible without requiring a financing equilibrium.

If the subscription settles on 29 June, the company pays all 100 at maturity.
If it settles on 2 July, the company first pays its available 60 at maturity,
records the unpaid 40 on 1 July, and pays the remaining 40 on 3 July. Both
paths finish with cash and debt equal to zero, inventory of 100 and equity of
100 after the interest expense. Only the second path has a missed-payment
history. The signed journal convention records debt as a credit balance;
the following table displays its positive carrying amount.

\begin{center}
\begin{tabular}{lrrrr}
\toprule
Late-funding path & Cash & Debt & Unpaid claim & Past miss \\
\midrule
Opening & 60 & 80 & 100 & 0 \\
Maturity, before payment & 60 & 100 & 100 & 0 \\
Partial payment, 30 June & 0 & 40 & 40 & 0 \\
Observation, 1 July & 0 & 40 & 40 & 40 \\
Share subscription, 2 July & 40 & 40 & 40 & 40 \\
Final payment, 3 July & 0 & 0 & 0 & 40 \\
\bottomrule
\end{tabular}
\end{center}

The historical amount in the last column is the amount recorded at the
missed-payment observation. It is not the current unpaid balance, and payment
does not reduce it. Neither that amount nor the missed flag is, by itself,
a finding that a particular legal event of default occurred.

\paragraph{Check the state and the change separately.}
The selected financing declaration uses the existing typed requirement
evaluator for 26 snapshot conditions. Face, remaining claim, recorded arrears,
historical missed amount and share counts must be nonnegative. The remaining
claim cannot exceed face. Before maturity the unpaid claim remains face while
the carrying amount may differ; after the selected maturity operation, carrying
debt equals the remaining claim. A recorded miss requires a matured promise,
an observation after its due date and a positive historical missed amount.
The recorded arrears equal the remaining claim in this selected post-observation
state. An unpaid state observed after the due date cannot pass without the
explicit missed-payment observation.

These conditions are useful but insufficient. After full payment, a snapshot
with a historical missed amount of 50 could satisfy the inequalities even
though the recorded amount was 40. Likewise, a nonnegative share count can
remain unchanged when a subscription should have increased it. The transition
declaration therefore binds the old state, typed event inputs and new state.
For a late payment of amount \(p\), let \(U\) be the remaining claim,
\(A\) the recorded arrears and \(H\) the historical missed amount. The
selected change requires
\[
  U_{\mathrm{after}}=U_{\mathrm{before}}-p,\qquad
  A_{\mathrm{after}}=A_{\mathrm{before}}-p,\qquad
  H_{\mathrm{after}}=H_{\mathrm{before}}.
\]
Cash falls by \(p\), the signed debt balance increases by \(p\), and the
missed flag remains true. Original face, due date and old shares are preserved.
A subscription has a different relation: new shares increase by the supplied
share quantity, while cash and capital change by quantity times issue price.
Money and shares have different dimensions, so a money-typed value cannot be
used as a share count merely because both are exact rationals.

The five declared transitions cover maturity, subscription, due-date payment,
missed-payment observation and late payment. They also preserve unrelated
memory, account balances and prior cash-flow records. Deliberately removing
the share-memory or historical-miss write can still yield a balanced proposal
from the generic engine. The selected transition relation rejects both. The
failed proposal remains available for diagnosis, while the accepted current
state and unexecuted command suffix retain the last valid prefix.

\paragraph{An internally consistent history still needs an opening premise.}
Opening assessment binds the selected method registry, schemas, reporting map,
account roles and typed financial state. A separate supplied history record
names the exact opening book and its debt-term fields, with nonblank provenance.
Missing or partial history is unresolved. A paid opening that erases the missed
flag and amount can be internally consistent, but it conflicts with the retained
history record and is refused. A caller who supplies the same false account
in both records can still satisfy this consistency test; source evidence must
establish the historical facts.

A valid intermediate snapshot is also not automatically a fiscal opening.
The existing statement engine requires the opening date to equal the fiscal
period start, an empty current-period cash-flow sequence and zero temporary
account balances. The first verification attempt incorrectly expected an
intermediate funded state to qualify as an opening. Its refusal exposed the
test error. The corrected comparison checks snapshot and supplied-history
consistency separately while preserving the fiscal-opening restrictions.
Closing and carrying the financing history require an explicit further step.

\paragraph{Keep current contractual findings apart from economic history.}
The conditional remedy examples introduce a separate record for an identified
hypothetical series. It contains the selected clause finding, its observation
date, the acceleration finding and the evidence identifier. The applicability
of the selected clause is a supplied premise. A known unpaid principal amount
must agree with the outstanding claim in the economic book; an unknown
applicability finding remains unresolved rather than becoming false.

The retained conditional reference distinguishes seven kinds of question:
clause satisfaction, authorized notice, new acceleration, remaining acceleration
after rescission, protected-change consent, majority waiver and a finding about
another series. Explicit typed predicates reproduce its 58 hypothetical
classifications, including unknown facts and unresolved boundary-day order.
Only three of these routes are composed into state updates here: a clause
observation, a new acceleration assessment and a rescission assessment. The
classification table is not a complete contract execution system.

In the composed example, the funded but still unpaid company first receives
a supplied conditional principal-clause finding. That observation changes no
cash and does not itself accelerate the obligation. A separate, supported
declaration establishes acceleration. The existing typed settlement then pays
the 40. A fresh paid-clause observation becomes unsatisfied, but a lack of a
new acceleration trigger leaves the earlier acceleration finding intact.
The selected rescission assessment changes that state only when all of its
supplied conditions are resolved as required. The historical missed amount
remains 40 throughout. These operations record conditional contractual status;
priority, penalties, accelerated cash amounts and insolvency consequences need
their own economic and accounting methods.

The initial contract-status record and its schema extension are explicit
synthetic premises. They do not establish that the retained issuer description
governs this promise. The source-dependent conditions retain their original
scope and unresolved operative-instrument dependencies. Known automatic and
declaration routes also keep their different information requirements: the
automatic branch does not require a notice finding that it does not use.

\paragraph{A matching date does not establish that a finding is current.}
The first composed candidate required a clause finding from the same date as
the proposed acceleration. Its original trace passed, but a further comparison
found a failure: record the unpaid clause, pay the full claim later in the same
dated sequence, then request acceleration from the old finding. The date still
matches although the economic state has changed. That candidate incorrectly
accepts the stale finding relative to its own current-clause requirement.

The successor records a digest of a fixed economic-state projection with each
resolved clause observation. The projection contains the book's identity and
period, account definitions, balances, cash account, cash-flow history,
non-contract memory and its types. Conditional-status memory and seen-event
bookkeeping are excluded, so recording a remedy does not invalidate itself.
The selected schema and types are checked separately. A new acceleration
requires both the date and economic-state digest to match the stored observation.

The same-day payment now makes the stale request unresolved before any accepted
write. A new observation of the paid claim supports the consistent assessment
that there is no new trigger. Balanced account changes, changes to unrelated
economic memory and revisions to cash-flow evidence also invalidate the old
observation. This conservative rule avoids guessing which economic changes
matter. Matching digests establish equality of the supplied record, not truth
of the evidence or persistence of external legal facts. The relevant external
findings still have to be supplied for the event.

The selected financing contract reproduces four complete financing paths on
two account charts, including their statements. The repaired remedy composition
reproduces eight complete transitions, with separate same-day failure and repair
evidence. These checks support the stated finite declarations and preserve the
rejected predecessor. Wider financing decisions, carry, changed rights, actual
series admission and independent review remain necessary for the full design.

\begingroup\raggedright
The independent settlement and conditional-source records are retained under
\path{data/accounting-master/gap-program-20261026/}.
The selected financial conditions and opening checks are under
\path{data/accounting-master/financing-opening-invariants-20261026/}.
The date-only failure and its preserved comparator are under
\path{data/accounting-master/financing-remedy-binding-20261026/};
the checked successor is under
\path{data/accounting-master/financing-remedy-state-binding-20261026/}.
\par\endgroup

